以下在模素数 $p$ 下运算，输入项在 $[0,p)$。已知数列满足阶数不超过 $k$ 的线性递推时，取前 $2k$ 项可恢复最小递推；任意有限前缀得到的规律不保证能外推。

返回 $c=\{1,c_1,\ldots,c_L\}$，$L=c.size()-1$ 为递推阶数，满足
$$\begin{aligned}
\sum_{j=0}^{L}a_{i-j}c_j&=0\qquad(i\ge L),\\
a_i&=-\sum_{j=1}^{L}c_j a_{i-j}.
\end{aligned}$$
如果删掉最后的取负和插入 $1$，返回的 $L$ 个系数 $v_j$ 满足 $a_i=\sum_{j=0}^{L-1}v_j a_{i-j-1}$。

接 5.5 时，\texttt{LinearRec} 的 \texttt{trans} 为 $\{-c_1,\ldots,-c_L\}$；另两份的系数数组为 $\{0,-c_1,\ldots,-c_L\}$。负数均规范到 $[0,p)$，初始值取 $a_0,\ldots,a_{L-1}$。$L=0$ 时按零阶递推直接返回 $0$。

\inputminted{cpp}{sections/Polynomial/assets/Math/Berlekamp-Massey.cpp}

向量序列可用固定随机行向量投影为标量序列；矩阵幂序列取 $s_i=u^T A^i v$。投影后的最小多项式仅保证整除原最小多项式，随机化成功时才相等；复用到整个向量时须验证向量递推，失败重新投影。

\paragraph*{优化矩阵快速幂DP}

	设 $f_i=A^i f_0$ 有 $n$ 维，目标为固定线性组合 $s_i=u^T f_i$。求出前 $2n$ 个标量项，跑 Berlekamp-Massey，再调用快速线性递推。

\paragraph*{求矩阵最小多项式}

	$n\times n$ 矩阵 $A$ 的最小多项式是满足 $q(A)=0$ 的最低次\textbf{首一}多项式，次数不超过 $n$。对 $s_i=u^TA^iv$ 跑 BM，返回的 $c$ 对应
	$$q(x)=x^L+c_1x^{L-1}+\cdots+c_L.$$
	注意 $c$ 按\textbf{高次到低次}存放，与前面多项式运算的顺序相反。只需递推 $A^iv$，不用保存 $A^i$。可多次投影取所得 $q$ 的最小公倍式；确定性验证须检查 $q(A)=0$。

\paragraph*{求稀疏矩阵的行列式}

	本方法要求 $A$ \textbf{非奇异}。取对角元非零的随机对角阵 $B$，求 $AB$ 的投影最小多项式 $q$。在足够大的域上，随机化成功时 $q$ 等于特征多项式；只有当 $\deg q=n$ 时才按
	$$\det A=(-1)^n c_n/\det B$$
	计算，次数不足则重新随机。次数不足不能直接判定 $A$ 奇异。奇异情形需另行判定并返回 $0$，不能保证其最小多项式等于特征多项式。

\paragraph*{求稀疏矩阵的秩}

	设 $A$ 为 $n\times m$，随机取对角元非零的 $n\times n$ 对角阵 $P$、$m\times m$ 对角阵 $Q$，构造 $M=QA^TPAQ$，其尺寸为 $m\times m$。
	在足够大的素数域（$p>m$）上，预处理及投影成功时，$M$ 的最小多项式去掉所有 $x$ 因子后的次数就是 $\operatorname{rank}A$。这是随机算法，可独立重复并取最大值以降低失败概率。计算 $Mv$ 时从右到左依次作用 $Q,A,P,A^T,Q$，不用显式构造 $M$。

\paragraph*{解稀疏方程组}

	$Ax=b$，要求 $A$ 在模 $p$ 下\textbf{满秩}，尺寸为 $n\times n$，$b,x$ 为 $n\times1$ 列向量。若 $\{A^ib\}$ 的递推系数为 $\{c_0,\ldots,c_L\}$，其中 $c_0=1$、$L\le n$、$c_L\ne0$，则系数数量为 $L+1\le n+1$，且
	$$A^{-1}b=-\frac1{c_L}\sum_{i=0}^{L-1}c_{L-1-i}A^ib.$$
	每次尝试用新投影生成前 $2n$ 个标量项，求 BM 后再次递推 $A^ib$ 累加答案；仅当验算 $Ax=b$ 成功时返回，失败重试。每次尝试时间 $O(n\,\operatorname{nnz}(A)+n^2)$，额外空间 $O(n)$，其中 $\operatorname{nnz}(A)$ 为非零元数量。可选第三参数 \texttt{seed} 用于复现，默认随机初始化。
	
	\inputminted{cpp}{sections/Polynomial/assets/Math/解稀疏方程组.cpp}
